%==============================================================================
% Appendix B: Native local action and DBR computational checks
%==============================================================================
\section{Native local action and DBR computational checks}
\label{app:native}

\subsection{Reduced action ingredients}
\label{app:native:action}

Under the homogeneous background ansatz (Section~\ref{sec:setup:perturbation}),
the first-order EC+NY action reduces to a finite-dimensional reduced Lagrangian.
The independent variables are the scale factors $q_{a}(t)$, the lapse $N(t)$, the
homogeneous connection components, and the torsion ansatz scalars $(\eta,V)$.
The perturbation analysis uses the quadratic action obtained by block-decomposing
this reduced system into scalar / spin-1 / spin-2 sectors.

\subsection{Symbolic Dirac--Bergmann engine}
\label{app:native:engine}

The constraint analysis was performed with a symbolic computation engine.  The
procedure follows the standard Dirac--Bergmann
algorithm~\cite{ref10,ref11,ref12}:
\begin{enumerate}
  \item Legendre transform and extraction of the primary constraints.
  \item Generation of the secondary constraints by the consistency conditions
        (iterated to closure).
  \item Classification of first/second class, and construction of the Dirac
        bracket for the second-class pairs.
  \item Presentation of the reduced spectrum by Routh-type elimination after
        closure (the DBR order).
\end{enumerate}

The decisions of the engine do not fix in advance the expected values of any
particular topology.  In the implementation, each topology local builder is
passed through the scalar-axis DBR closure, and the activation entry is read
from primary preservation, secondary generation, the constraint Poisson matrix
rank, and the first/second-class split.  The indeterminate-symbol relation
between the structure constants and the class-A matrix is checked separately as
a formal identity.  The physical role of that map as DBR activation is checked
by the axis-by-axis closure audit.

\paragraph{Correspondence with the proof skeleton of
Section~\ref{sec:rankdefect:mechanism}.}
Write $\Aop_{\rm DBR}$ for the secondary activation map for the scalar-axis
perturbations $\bm{\beta}\in\mathbb{R}^{3}$.  In the class-A Milnor frame, the
formal structure-map identity gives
\begin{equation}
  \Aop_{\rm form} = n .
\end{equation}
This identity is a notation and structure-map check; it is not used as a
substitute for DBR closure.  The witness
\code{rank_defect_activation_witness.py} calls the bundled DBR
constraint-closure engine (\code{dppu.dirac}) and computes the entries of
$\Aop_{\rm DBR}$ from DBR closure for the four topologies and three scalar
axes.  The resulting engine-derived activation matrix satisfies
$\rank(\Aop_{\rm DBR})=\rank(n)$ and
$\dim\ker(\Aop_{\rm DBR})=\dzero$.

\subsection{Rank verification conventions}
\label{app:native:rank}

For the rank decisions of the constraint matrices and the kinetic Hessians we
impose the following redundancies:
\begin{itemize}
  \item \textbf{Two routes}: agreement of the symbolic rank and the numerical
        rank is required.
  \item \textbf{Two generic points}: the numerical rank is evaluated at two
        generic background points, and agreement is required (excluding an
        accidental single-point degeneracy).
  \item \textbf{Mod-$p$ check}: for large symbolic matrices, a rank computation
        at a prime modulus is used in parallel to exclude a spurious rank from
        cancellation/round-off.
\end{itemize}

\subsection{Frame-invariance check protocol}
\label{app:native:frame}

The frame invariance of Section~\ref{sec:rankdefect:frame} is based first on the
analytic argument that $\rank(n)$ and $\dim\ker n$ are preserved under the
congruence $n\mapsto RnR^{\mathsf{T}}$.  The numerical verification reinforces it
by two routes: (a) random SO(3) rotations ($\dim\ker n$ preserved in all 20
rotations), and (b) reconstruction of $n$ from the rotated structure constants
(tensor transformation) and recomputation of the kernel dimension.  The script
\code{rank_defect_activation_witness.py} also confirms
$\rank(RnR^{\mathsf{T}})=\rank(n)$ for rational orthogonal matrices $R$.

\subsection{Comparison procedure under a single protocol}
\label{app:native:protocol}

The comparison of the four topologies was carried out under one and the same
engine, the same sector decomposition, and the same rank-verification conventions.
$\Sthree$ doubles as a cross-check of the method, passing the null check of the
scalar-axis projection (all axes zero for the isotropic $\Sthree$).

\subsection{Reproducibility notes}
\label{app:native:repro}

The main formulas and decisions in the paper are reproducible by the independent
witness scripts under \code{script/scripts/paper07/}.

\begin{table}[H]
\centering
\begingroup
\renewcommand{\arraystretch}{1.3}
\centering
\resizebox{\linewidth}{!}{%
\begin{tabular}{ll}
\toprule
Witness & Corresponding sections \\
\midrule
\code{rank_defect_activation_witness.py} &
\S\ref{sec:rankdefect:mechanism}--\ref{sec:rankdefect:frame},
\ref{app:native:engine}--\ref{app:native:frame}; *1 \\
\code{spectral_completion_witness.py} &
\S\ref{sec:spectral}, \ref{app:spectral} \\
\code{torsion_auxiliary_field_promotion_witness.py} &
\S\ref{sec:torsion}, \ref{app:torsion} \\
\code{finite_fate_branch_witness.py} &
\S\ref{sec:finite}, \ref{app:finite} \\
\code{harmonic_kernel_correspondence.py} &
\S\ref{sec:harmonic}, \ref{app:harmonic} \\
\code{route_comparison_s3_branch_a_witness.py} &
\S\ref{sec:dbr:control}, \ref{app:route} \\
\code{s3_normalization_factor_witness.py} &
\S\ref{sec:setup:structure}, \ref{app:conventions:s3} \\
\bottomrule
\end{tabular}%
}
\endgroup
\caption{Independent witness scripts and the sections they reproduce. (*1: formal identity as notation
check, DBR closure audit as primary evidence)}
\label{tab:witness}
\end{table}
